\section{Limitations, failure modes, and counterexamples}
\label{sec:limitations}

The preceding sections establish the BayesBreak framework as a modular exact-Bayesian
segmentation pipeline with explicit theoretical guarantees on conjugate-block inference
(Theorem~\ref{thm:ef-integral}, Proposition~\ref{prop:bb-complexity}), posterior-odds
stability under non-conjugate approximation (Assumption~\ref{ass:uniform-block-error},
Proposition~\ref{prop:stability}, Proposition~\ref{prop:uniform-bounds}), latent-group
identifiability up to label permutation (Proposition~\ref{prop:latent-identifiability},
Example~\ref{ex:label-switch-counterexample}), and joint MAP correctness
(Theorem~\ref{thm:map-correctness}). This section collects the corresponding limitations
and named failure modes, so that downstream users can reason about applicability
before running the framework on a new problem class.

\paragraph{Computational regime.}
The DP layer is $\Theta(k_{\max} n^2)$ time and at least $\Theta(k_{\max} n)$ memory
(Proposition~\ref{prop:bb-complexity}). The framework is therefore appropriate for series
lengths $n$ in the thousands and modest segment ceilings $k_{\max}$; it is \emph{not} an
online detector and is not competitive with linear-time penalized methods (e.g.\ PELT) at
$n\gtrsim 10^5$ unless block precomputation is dramatically sparsified. The DP itself
cannot be replaced with a coarser approximation without losing the exact posterior
guarantees that motivate the framework; the right substitution for very large $n$ is a
sliding-window decomposition (planned, not implemented here).

\paragraph{Block-model misspecification.}
All guarantees are conditional on the chosen exponential-family block model. If the true
within-segment distribution is not in the chosen family (e.g.\ heavy-tailed observations
under a Gaussian block, or zero-inflation under a Poisson block), block evidences become
biased and the posterior over $k$ inherits that bias. The framework provides
posterior-predictive diagnostics for this case (Section~\ref{sec:families-predictive}),
but does \emph{not} adapt the block family automatically. A practical recipe is to
compare the implemented block families on the same data (Table~\ref{tab:single_quant})
and reserve the chosen specification for the body of the paper.

\paragraph{Non-conjugate approximation regime.}
The posterior-odds bound of Proposition~\ref{prop:stability} is linear in $k_{\max}$, so
non-conjugate analyses with $k_{\max}\gg 30$ must verify a correspondingly smaller $\varepsilon$
to keep $(k+k')\varepsilon$ small. Routine-by-routine $\varepsilon$ rates are in
Proposition~\ref{prop:uniform-bounds}; the EP row of Table~\ref{tab:nonconj_tradeoff} is the
canonical failure mode (EP need not satisfy a uniform $\varepsilon$ when its iteration fails to
converge).

\paragraph{Latent-group identifiability.}
Proposition~\ref{prop:latent-identifiability} identifies the latent-group template
mixture up to permutation of the group labels;
Example~\ref{ex:label-switch-counterexample} makes the indeterminacy explicit at the
two-group, two-subject scale. Any label-level reporting therefore requires a
deterministic anchoring convention (e.g.\ ordering templates by $k_g$ or by lexicographic
order on $t^{(g)}$). The EM algorithm of Section~\ref{sec:latent-em} is not guaranteed to
recover the global maximizer of the finite template-mixture objective; the
recommendation of $R\ge 5$ restarts in Section~\ref{sec:latent-em} is a defense against
this but is not a guarantee.

\paragraph{Boundary semantics.}
Boundary-event marginals $p(b_i=1\mid y,k)$ and the joint MAP segmentation
$\widehat t^{\mathrm{MAP}}(k)$ are different posterior summaries (Remark~\ref{rem:marg-vs-joint})
and must not be conflated. Report both whenever boundary semantics matter for the downstream task.

\paragraph{Partition-prior sensitivity.}
The posterior over segment counts $p(k\mid y)$ is mildly sensitive to the prior $p(k)$
and substantially sensitive to the length factor $g(\Delta_x(i,j))$ when the design grid is
irregular. The framework supplies design-aware partition priors
(Section~\ref{sec:irregular}); the well-log results table (Table~\ref{tab:real_welllog},
planned) is the natural place to quantify this sensitivity for users. Until that
ablation is populated, users are advised to compare an index-uniform prior to at least
one length-aware alternative and report any disagreement explicitly.

\paragraph{Out-of-scope settings.}
The framework does \emph{not} address: (i) streaming or online detection (use
\citet{adamsmackay2007bocpd}); (ii) within-segment trends or change-in-slope models (use
piecewise-polynomial regression, e.g.\ \citet{ruggieri2014bcpvs}); (iii) variable-dimensional
posterior simulation over $k$ with explicit posterior samples of $t$
(use~\citet{green1995rjMCMC,fearnhead2006exact}); (iv) network or spatial segmentation on
non-totally-ordered designs. The motivating reason to use BayesBreak is exact
sum-product inference under \emph{ordered} designs and \emph{piecewise-constant} block
means; the moment a problem falls outside that scope, a different framework is more
appropriate.

\paragraph{Scope of the empirical evaluation in this paper.}
The synthetic suite and the four real-data figures are the implemented demonstrations of the
framework; conclusions in this paper are restricted to claims that the implemented framework
recovers what it was designed to recover. Broader competitiveness claims require the planned
external comparators (PELT, wild binary segmentation, SMUCE, RJMCMC, Fearnhead's exact DP) and
the partition-prior / pooling / likelihood ablations populated by the finalized benchmarking
pipeline.
